\chapter{Episiotomy}

\section*{Overview}
An episiotomy is a surgical cut in the perineum, the area between the vaginal opening and the anus, made during vaginal birth. It widens the vaginal opening. It is not recommended as a routine step in spontaneous vaginal birth. A clinician may recommend one when there is a specific need to help the baby be born, such as an assisted birth with forceps or a suction cup, or concern that the baby needs to be delivered quickly.

\section*{Alternative Names}
Episiotomy; perineotomy. A mediolateral episiotomy is angled away from the vaginal opening; a midline episiotomy is made straight down toward the anus.

\section*{Principle}
The cut enlarges the vaginal opening to help with birth in selected situations. It does not prevent all perineal tears and can sometimes extend into a deeper tear. The clinician examines the vagina and perineum after birth to identify and repair the incision and any additional injury.

\section*{History}
Episiotomy was once used routinely in many settings. Evidence and modern guidelines support a more selective approach because routine use does not reliably prevent serious tears and exposes more people to pain and wound complications.

\section*{Why It Is Done}
An episiotomy may be considered when there is a clinical reason to assist vaginal birth, for example during some forceps or vacuum-assisted births, or when the baby needs to be born quickly because of concern about the baby's condition. It is not done simply to prevent tearing in an uncomplicated spontaneous birth. The clinician discusses the reason when circumstances allow; an urgent situation may limit time for discussion.

\section*{Role in Care}
Episiotomy is an incision made to assist vaginal birth, not a diagnostic test and not a routine part of birth. Whether it is appropriate depends on what is happening during labor, the baby's and birthing person's condition, and the available options.

\section*{How It Is Performed}
If time and circumstances allow, the clinician explains the reason and obtains consent, then provides effective pain relief, often by using an existing epidural or injecting local anesthetic. During the pushing stage, the clinician makes the incision in the perineum. NICE recommends a mediolateral incision, usually angled to the right, with an angle of 45 to 60 degrees from the vertical at the time it is made. Technique and guidance may vary by country. In an emergency because of acute concern for the baby, the incision may need to be made without delay. After the baby and placenta are delivered, the clinician examines the area and repairs the incision with dissolvable stitches, using additional expertise if a deeper tear is found.

\section*{Preparation}
There is usually no separate preparation because the decision is made during birth. The birth team explains the reason and pain relief when there is time. Tell the team about allergies, bleeding problems, medicines that affect bleeding, or concerns about previous perineal injury if you can do so before labor.

\section*{Contraindications and Precautions}
Episiotomy is generally avoided when there is no clear clinical need. It is not automatically required for an assisted birth; the clinician considers the circumstances and alternatives. A previous third- or fourth-degree tear or other perineal injury may affect planning, but does not by itself determine whether an episiotomy is needed in a later birth. If a tear extends into the anal sphincter or rectum, it requires careful assessment, repair by a clinician with appropriate expertise, and follow-up.

\section*{Risks}
Pain and soreness, bleeding, bruising or a blood collection (hematoma), infection, and delayed healing or separation of the wound can occur. Pain with sex may persist for some people. The cut can extend into a more serious tear involving the anal sphincter or rectum. The risk and type of injury vary; an episiotomy does not guarantee that a severe tear will be avoided.

\section*{Aftercare and Recovery}
Soreness is common, especially when sitting or walking, and usually improves over the first few weeks. Dissolvable stitches do not need removal. Keep the area clean with plain water, change sanitary pads regularly, and wash hands before and after using the toilet or changing pads. Follow the maternity team's advice about pain relief, cold packs, bowel care, and pelvic floor exercises. Avoid sex or activities that cause pain until the wound has healed and you feel ready; there is no single timetable that suits everyone.

\section*{Outcome and Follow-up}
The incision is repaired after birth and usually heals over several weeks. The clinician checks for bleeding, infection, wound separation, and deeper injury. Follow-up is especially important if the cut extends into the anal sphincter, if healing is difficult, or if pain or problems controlling wind or stool continue.

\section*{Expected Results}
The goal is to make more room for birth when clinically needed. Whether an episiotomy helps depends on the situation, and it does not ensure an uncomplicated birth or prevent every tear. Most incisions heal without a lasting problem, although discomfort can take time to settle.

\section*{Follow-up}
Contact the midwife, obstetric clinician, or primary care clinician if pain is not improving, the wound appears open, or there are concerns about healing. Persistent pain during sex, ongoing perineal pain, or difficulty controlling wind or stool deserves assessment and support. People with a repaired third- or fourth-degree tear should follow the specific specialist follow-up plan.

\section*{When to Seek Medical Care}
Seek prompt medical advice for increasing pain, a fever, foul-smelling discharge, new redness or swelling, pus, wound separation, or heavy bleeding. Contact a clinician if you have trouble controlling wind or stool, cannot pass urine, or have new severe pain. Seek urgent care for very heavy bleeding, fainting, or rapidly worsening symptoms.

\section*{References}
\begin{enumerate}
  \item National Institute for Health and Care Excellence. Intrapartum care. https://www.nice.org.uk/guidance/ng235/chapter/Recommendations
  \item World Health Organization. WHO recommendations on maternal health: guidelines approved by the WHO Guidelines Review Committee. https://iris.who.int/bitstream/handle/10665/380666/9789240080591-eng.pdf?sequence=1
  \item Royal College of Obstetricians and Gynaecologists. Episiotomy. https://www.rcog.org.uk/for-the-public/perineal-tears-and-episiotomies-in-childbirth/episiotomy/
  \item MedlinePlus. Episiotomy - aftercare. https://medlineplus.gov/ency/patientinstructions/000483.htm
\end{enumerate}

\clearpage
